% GENERATED FILE. Edit canonical lessons, then run npm run generate:books.
% canonical-source-hash: fnv1a64:ea6141c208b3b997
% canonical-lessons: LA-C457-insula, LA-C457-cenaculum, LA-C457-contignatio, LA-C457-tignum, LA-C457-trabs

\chapter{Block of flats, Upper room, Storey, Timber, Beam}
\label{ch:la-block-of-flats-upper-room-storey-timber-beam}
\hlchaptermodality{457}

\begin{chapteropening}
\textbf{By the end of this chapter:} \emph{I can describe a block of flats: the block itself, an upper room, a storey, timbers and beams.}

Complete the last lesson of chapter 457: I can describe a block of flats: the block itself, an upper room, a storey, timbers and beams.
\end{chapteropening}

\section[īnsula, īnsulae]{īnsula, īnsulae — a block of flats, an apartment house; an island}

\label{lesson:LA-C457-insula}



\begin{quote}
Before the new one: say the Latin for to live in, then the Latin for an occupant.
\end{quote}

\subsection*{You'll want to know: īnsula, īnsulae}
\textbf{īnsula, īnsulae} — \textquotedblleft{}a block of flats, an apartment house; an island\textquotedblright{}.

\begin{quote}
In urbe multī pauperēs in īnsulīs altīs habitant.
\end{quote}

\begin{quote}
\emph{In the city many poor people live in tall blocks of flats.}
\end{quote}

First declension, feminine: \emph{īnsula, īnsulae f.} First an island, then a block of rented flats standing alone like an island among streets. English \emph{isle} and \emph{insulate} come from it.

\begin{cousinweb}
The origin of \emph{īnsula} is not settled: Lewis \& Short offer a doubtful analysis as \emph{in-} + \emph{sul-}, and EtymDB gives no older parent. English \emph{isle} comes from it through French, and \emph{insulate} and \emph{peninsula} from Latin.
\end{cousinweb}

\begin{grammarlens}[title={What you've built}]
One more word for telling, explaining and describing at B1.
\end{grammarlens}

\subsection*{Your turn}
\begin{itemize}
\raggedright
  \item \emph{Say it:} \emph{īnsula, īnsulae}
  \item \emph{Say it:} \emph{īnsula, īnsulae}, once more
  \item \emph{Say it:} a sentence of your own with \emph{īnsula, īnsulae}
  \item \emph{Recall:} say the Latin for to live in, then the Latin for an occupant, then say \emph{īnsula, īnsulae} again
\end{itemize}

\begin{tcolorbox}[breakable,colback=teal!4,colframe=teal!35!black,title={Before you move on}]
What does īnsula, īnsulae mean? (\textquotedblleft{}A block of flats, an apartment house; an island\textquotedblright{}.) Say it once more.
\end{tcolorbox}
\section[cēnāculum, cēnāculī]{cēnāculum, cēnāculī — an upper room, a garret, a flat on an upper floor}

\label{lesson:LA-C457-cenaculum}



\begin{quote}
Before the new one: say the Latin for an occupant, then the Latin for a block of flats.
\end{quote}

\subsection*{You'll want to know: cēnāculum, cēnāculī}
\textbf{cēnāculum, cēnāculī} — \textquotedblleft{}an upper room, a garret, a flat on an upper floor\textquotedblright{}.

\begin{quote}
Poēta pauper in cēnāculō sub tēgulīs habitat.
\end{quote}

\begin{quote}
\emph{The poor poet lives in a garret under the roof tiles.}
\end{quote}

Second declension, neuter: \emph{cēnāculum, cēnāculī n.} First a dining room, which in old houses was upstairs, then any upper room.

\begin{cousinweb}
\emph{cēnāculum} is built on \emph{cēna} 'dinner' with the suffix \emph{-culum} 'place for' (EtymDB; Lewis \& Short): a dining room, which was upstairs, and so an upper room.
\end{cousinweb}

\begin{grammarlens}[title={What you've built}]
One more word for telling, explaining and describing at B1.
\end{grammarlens}

\subsection*{Your turn}
\begin{itemize}
\raggedright
  \item \emph{Say it:} \emph{cēnāculum, cēnāculī}
  \item \emph{Say it:} \emph{cēnāculum, cēnāculī}, once more
  \item \emph{Say it:} a sentence of your own with \emph{cēnāculum, cēnāculī}
  \item \emph{Recall:} say the Latin for an occupant, then the Latin for a block of flats, then say \emph{cēnāculum, cēnāculī} again
\end{itemize}

\begin{tcolorbox}[breakable,colback=teal!4,colframe=teal!35!black,title={Before you move on}]
What does cēnāculum, cēnāculī mean? (\textquotedblleft{}An upper room, a garret, a flat on an upper floor\textquotedblright{}.) Say it once more.
\end{tcolorbox}
\section[contignātiō, contignātiōnis]{contignātiō, contignātiōnis — a storey, a floor (of beams and boards)}

\label{lesson:LA-C457-contignatio}



\begin{quote}
Before the new one: say the Latin for a block of flats, then the Latin for an upper room.
\end{quote}

\subsection*{You'll want to know: contignātiō, contignātiōnis}
\textbf{contignātiō, contignātiōnis} — \textquotedblleft{}a storey, a floor (of beams and boards)\textquotedblright{}.

\begin{quote}
Īnsula alta est: sex contignātiōnēs habet.
\end{quote}

\begin{quote}
\emph{The block of flats is tall: it has six storeys.}
\end{quote}

Third declension, feminine: \emph{contignātiō, contignātiōnis f.} A floor laid on joists, so one storey of a building.

\begin{cousinweb}
\emph{contignātiō} is built on \emph{contignō} 'to join with beams', from \emph{con-} and \emph{tignum} 'beam, timber' (EtymDB; Lewis \& Short).
\end{cousinweb}

\begin{grammarlens}[title={What you've built}]
One more word for telling, explaining and describing at B1.
\end{grammarlens}

\subsection*{Your turn}
\begin{itemize}
\raggedright
  \item \emph{Say it:} \emph{contignātiō, contignātiōnis}
  \item \emph{Say it:} \emph{contignātiō, contignātiōnis}, once more
  \item \emph{Say it:} a sentence of your own with \emph{contignātiō, contignātiōnis}
  \item \emph{Recall:} say the Latin for a block of flats, then the Latin for an upper room, then say \emph{contignātiō, contignātiōnis} again
\end{itemize}

\begin{tcolorbox}[breakable,colback=teal!4,colframe=teal!35!black,title={Before you move on}]
What does contignātiō, contignātiōnis mean? (\textquotedblleft{}A storey, a floor (of beams and boards)\textquotedblright{}.) Say it once more.
\end{tcolorbox}
\section[tignum, tignī]{tignum, tignī — a timber, a log for building; building material}

\label{lesson:LA-C457-tignum}



\begin{quote}
Before the new one: say the Latin for an upper room, then the Latin for a storey.
\end{quote}

\subsection*{You'll want to know: tignum, tignī}
\textbf{tignum, tignī} — \textquotedblleft{}a timber, a log for building; building material\textquotedblright{}.

\begin{quote}
Fabrī tigna in plaustrō ad īnsulam novam vehunt.
\end{quote}

\begin{quote}
\emph{The workmen carry timbers on a wagon to the new block of flats.}
\end{quote}

Second declension, neuter: \emph{tignum, tignī n.} Wood for building, cut and ready.

\begin{cousinweb}
\emph{tignum} is inherited from Proto-Italic \emph{tegnom} (EtymDB), which EtymDB traces to a PIE word for a pole or stake, with English \emph{stake} as a cousin. Lewis \& Short connect it instead with the root of \emph{texō} 'to build', so the deeper source is not settled.
\end{cousinweb}

\begin{grammarlens}[title={What you've built}]
One more word for telling, explaining and describing at B1.
\end{grammarlens}

\subsection*{Your turn}
\begin{itemize}
\raggedright
  \item \emph{Say it:} \emph{tignum, tignī}
  \item \emph{Say it:} \emph{tignum, tignī}, once more
  \item \emph{Say it:} a sentence of your own with \emph{tignum, tignī}
  \item \emph{Recall:} say the Latin for an upper room, then the Latin for a storey, then say \emph{tignum, tignī} again
\end{itemize}

\begin{tcolorbox}[breakable,colback=teal!4,colframe=teal!35!black,title={Before you move on}]
What does tignum, tignī mean? (\textquotedblleft{}A timber, a log for building; building material\textquotedblright{}.) Say it once more.
\end{tcolorbox}
\section[trabs, trabis]{trabs, trabis — a beam, a timber; (in poetry) a ship}

\label{lesson:LA-C457-trabs}



\begin{quote}
Before the new one: say the Latin for a storey, then the Latin for a timber.
\end{quote}

\subsection*{You'll want to know: trabs, trabis}
\textbf{trabs, trabis} — \textquotedblleft{}a beam, a timber; (in poetry) a ship\textquotedblright{}.

\begin{quote}
Magnae trabēs in tēctō templī sunt.
\end{quote}

\begin{quote}
\emph{There are great beams in the roof of the temple.}
\end{quote}

Third declension, feminine: \emph{trabs, trabis f.} Bigger than a \emph{tignum}. English \emph{architrave} comes, through Italian \emph{trave}, from it.

\begin{cousinweb}
\emph{trabs} is inherited from PIE \emph{treb-} 'wooden beam' (EtymDB), and EtymDB lists English \emph{thorp} 'village' as a cousin. English \emph{architrave} comes from it through Italian \emph{trave}.
\end{cousinweb}

\begin{grammarlens}[title={What you've built}]
One more word for telling, explaining and describing at B1.
\end{grammarlens}

\subsection*{Your turn}
\begin{itemize}
\raggedright
  \item \emph{Say it:} \emph{trabs, trabis}
  \item \emph{Say it:} \emph{trabs, trabis}, once more
  \item \emph{Say it:} a sentence of your own with \emph{trabs, trabis}
  \item \emph{Recall:} say the Latin for a storey, then the Latin for a timber, then say \emph{trabs, trabis} again
\end{itemize}

\begin{tcolorbox}[breakable,colback=teal!4,colframe=teal!35!black,title={Before you move on}]
What does trabs, trabis mean? (\textquotedblleft{}A beam, a timber; (in poetry) a ship\textquotedblright{}.) Say it once more.
\end{tcolorbox}
